\chapter{自然语言处理：应用}

预训练向量或 BERT 编码器提供 $w\mapsto\bm{v}$ 或序列 $\mapsto\mathsf{H}$。
下游把可变长文本映到固定标签：情感、句对关系、词元标签或问答跨度。

\section{情感分析及数据集}

情感分析是文本分类：序列 $\bm{w}_{1:T}$ 映到极性 $y\in\{0,1\}$。
数据为 IMDb 影评，正负标签。

\subsection{读取数据集}

每个样本是评论文本与 $y\in\{0,1\}$。

\subsection{预处理数据集}

词元为词，低频过滤后建词表。长度直方图显示 $T$ 变化大；截断或填充到固定 $L=500$。

\subsection{创建数据迭代器}

每个批量返回 $(\bm{X},\bm{y})$，$\bm{X}\in\mathbb{Z}^{B\times L}$ 为词元下标。

\subsection{整合代码}

加载函数返回训练/测试迭代器与词表。

\subsection{小结}

情感是
$\{w_t\}_{t=1}^{T}\mapsto y\in\{0,1\}$。
变长 $T$ 经截断填充变成固定 $L$。

\subsection{练习}

核验：哪些预处理超参缩短每个 epoch，以及如何把另一评论语料映到同一迭代器接口。

\section{情感分析：使用循环神经网络}

小数据集上用冻结的预训练词向量，再用双向 RNN 把序列压成一个向量。

\subsection{使用循环神经网络表示单个文本}

嵌入 $\bm{e}_t=\bm{E}[w_t]$（GloVe）。
双向 LSTM 最后一层的初始与最终时刻隐状态拼接为 $\bm{h}$，再
\[
\hat{\bm{y}}=\operatorname{softmax}(\bm{W}\bm{h}+\bm{b})\in\mathbb{R}^{2}.
\]

\subsection{加载预训练的词向量}

$\bm{E}$ 用 $100$ 维 GloVe 初始化，训练中不更新，使 $\bm{e}_t$ 固定。

\subsection{训练和评估模型}

最小化交叉熵。推理时把新句词元化、填充后取 $\hat y=\arg\max\hat{\bm{y}}$。

\subsection{小结}

\[
\bm{h}=[\overrightarrow{\bm{h}}_1;\overleftarrow{\bm{h}}_T],
\qquad
\hat{\bm{y}}=\operatorname{softmax}(\bm{W}\bm{h}+\bm{b}),
\]
其中 $\bm{e}_t$ 来自冻结 GloVe。

\subsection{练习}

核验：增大周期或改用 $300$ 维 GloVe 是否降低测试误差。

\section{情感分析：使用卷积神经网络}

把序列看成一维“图像”，用卷积抽 $n$ 元语法，再用时间最大汇聚得到定长向量。

\subsection{一维卷积}

核 $\bm{k}$ 在序列上滑动，输出
\[
y_i=\sum_{j=0}^{k-1}x_{i+j}k_j.
\]
多通道时等价于把嵌入维当作二维卷积的通道。

\subsection{最大时间汇聚层}

每个通道在时间维取 $\max$，得到与通道数相同的向量，允许各通道有效长度不同。

\subsection{textCNN模型}

输入形状视为宽 $n$、高 $1$、通道 $d$。
若干不同宽度的一维核抽局部特征；各通道时间最大汇聚后拼接；全连接加 Dropout 得到两类。

\subsubsection{定义模型}

两套嵌入：一套可训练，一套固定 GloVe；通道上拼接后再卷积。
例如核宽 $\{3,4,5\}$，每类 $100$ 个输出通道。

\subsubsection{加载预训练词向量}

可训练表与固定表都用 $100$ 维 GloVe 初始化，后者不更新。

\subsubsection{训练和评估模型}

交叉熵训练。推理同 RNN 节。

\subsection{小结}

一维卷积抽 $n$ 元，时间最大汇聚得
$\bm{z}=\operatorname{concat}_c\max_t(\mathsf{Y}_{c,t})$，再线性分类。

\subsection{练习}

核验：与双向 LSTM 在精度和吞吐上的差别，以及加上位置编码是否改变 $\hat y$。

\section{自然语言推断与数据集}

NLI 判断假设能否由前提推出。标签为蕴涵、矛盾、中性。

\subsection{自然语言推断}

前提 $A$、假设 $B$ 都是词元序列。输出
\[
y\in\{\text{蕴涵},\text{矛盾},\text{中性}\}.
\]

\subsection{斯坦福自然语言推断（SNLI）数据集}

约 $5\times 10^5$ 带标签英文句对。

\subsubsection{读取数据集}

抽取 $(A,B,y)$。打印可见三类在训练与测试上大致均衡。

\subsubsection{定义用于加载数据集的类}

长度统一为 $\mathrm{num\_steps}$：超长截断，不足填填充词元。
下标 $i$ 返回一对前提、假设与标签。

\subsubsection{整合代码}

词表只由训练集构建。批量里有两个输入张量，分别是前提与假设。

\subsection{小结}

NLI 是句对分类 $y=f(A,B)$。SNLI 提供三类均衡的监督。

\subsection{练习}

核验：用 NLI 模型给机器翻译候选打“是否蕴涵参考”的分数是否可行。

\section{自然语言推断：使用注意力}

可分解注意力：对齐、比较、聚合，无循环、无卷积。

\subsection{模型}

前提词向量 $\bm{a}_1,\ldots,\bm{a}_m$，假设 $\bm{b}_1,\ldots,\bm{b}_n$。

\subsubsection{注意（Attending）}

\[
e_{ij}=f(\bm{a}_i)^{\mathsf T}f(\bm{b}_j).
\]
软对齐
\begin{align*}
\bm{\beta}_i
&=
\sum_{j=1}^{n}
\frac{\exp(e_{ij})}{\sum_{k=1}^{n}\exp(e_{ik})}\,\bm{b}_j,
\\
\bm{\alpha}_j
&=
\sum_{i=1}^{m}
\frac{\exp(e_{ij})}{\sum_{k=1}^{m}\exp(e_{kj})}\,\bm{a}_i.
\end{align*}
$f$ 先把向量映到较低维再做内积，复杂度对 $m,n$ 是线性而非二次深层交互。

\subsubsection{比较}

\begin{align*}
\bm{v}_{A,i}&=g([\bm{a}_i,\bm{\beta}_i]),\quad i=1,\ldots,m,
\\
\bm{v}_{B,j}&=g([\bm{b}_j,\bm{\alpha}_j]),\quad j=1,\ldots,n.
\end{align*}

\subsubsection{聚合}

\[
\bm{v}_A=\sum_{i=1}^{m}\bm{v}_{A,i},\qquad
\bm{v}_B=\sum_{j=1}^{n}\bm{v}_{B,j},
\]
\[
\hat{\bm{y}}=h([\bm{v}_A,\bm{v}_B]).
\]
$h$ 输出三类 logits。

\subsubsection{整合代码}

注意、比较、聚合联合训练，参数即 $f,g,h$ 中的 MLP。

\subsection{训练和评估模型}

在 SNLI 上优化交叉熵。

\subsubsection{读取数据集}

批量 $256$，序列长 $50$。

\subsubsection{创建模型}

$f(\bm{a}_i)$ 与词向量 $100$ 维；$f,g$ 输出 $200$ 维。GloVe 初始化词表。

\subsubsection{训练和评估模型}

句对作为两个输入分到多卡。记录训练与测试准确率。

\subsubsection{使用模型}

新的 $(A,B)$ 词元化后取 $\arg\max\hat{\bm{y}}$。

\subsection{小结}

对齐 $e_{ij}=f(\bm{a}_i)^{\mathsf T}f(\bm{b}_j)$，
比较 $g([\bm{a}_i,\bm{\beta}_i])$，聚合后 $h([\bm{v}_A,\bm{v}_B])$ 得三类。

\subsection{练习}

核验：该分解无法对高阶跨句交互建模时误差如何表现，以及如何改成输出 $[0,1]$ 相似度。

\section{针对序列级和词元级应用微调BERT}

BERT 编码后只加浅层。序列级用 $\langle\mathrm{cls}\rangle$；词元级用每个位置的向量。
额外层从头学，编码器全部微调。

\subsection{单文本分类}

输入一条序列。$\langle\mathrm{cls}\rangle$ 向量经线性层得 $K$ 类，如情感或语法可接受性。

\subsection{文本对分类或回归}

句对打包进一条 BERT 序列，片段嵌入区分两句。
分类：NLI 三类。回归：语义相似度分数，例如 $[0,5]$ 上的实数，损失用平方误差。

\subsection{文本标注}

每个词元的 BERT 向量送到共享（或位置共享）的分类头，如词性。
$\langle\mathrm{cls}\rangle$、$\langle\mathrm{sep}\rangle$ 通常不预测标签。

\subsection{问答}

段落与问题拼成句对。在段落词元上预测答案起点与终点的两个分类器，
答案跨度为 $\arg\max_{s\leq e}P_{\mathrm{start}}(s)P_{\mathrm{end}}(e)$。

\subsection{小结}

序列级：$\hat{\bm{y}}=h(\bm{h}_{\langle\mathrm{cls}\rangle})$。
词元级：$\hat y_t=h(\bm{h}_t)$。问答再加起点终点两个头。

\subsection{练习}

核验：检索中如何用负采样与 BERT 打分，以及 BERT 作编码器时 MLM 头如何改回标准 LM。

\section{自然语言推断：微调BERT}

把 SNLI 当作文本对分类，在预训练 BERT 上加 MLP。

\subsection{加载预训练的BERT}

小模型便于演示；base 与原 BERT 基础规模相当。
加载词表与预训练参数。

\subsection{微调BERT的数据集}

前提与假设打包为 BERT 输入，片段下标区分两句。
超长时轮流删较长一句的末词元直到不超过 $\mathrm{max\_len}$。

\subsection{微调BERT}

$\langle\mathrm{cls}\rangle$ 表示经两层全连接得三类。
编码器与嵌入参与微调；仅与预训练 MLM/NSP 头相关、且未接入该分类器的参数不更新。

\subsection{小结}

\[
\hat{\bm{y}}
=
\operatorname{MLP}\bigl(\bm{h}_{\langle\mathrm{cls}\rangle}(A,B)\bigr)
\in\mathbb{R}^{3}.
\]
BERT 成为下游模型的子图，预训练任务头不再前向。

\subsection{练习}

核验：按两句长度比截断与轮流删末词元哪种更少切断前提或假设的关键词。
